\subsubsection{Plano de Testes}

\paragraph{Estratégia de Testes}

O plano de testes foi elaborado com base nos requisitos funcionais (Seção~3.1.1) e nos critérios de aceitação (Quadro~\ref{qua:criterios}), contemplando exclusivamente os módulos implementados no escopo do MVP. A estratégia adotada prioriza testes unitários como primeira camada de defesa, combinada com testes de integração para validação das camadas de infraestrutura e testes \textit{end-to-end} (E2E) para verificação do fluxo completo do sistema.

\begin{itemize}
    \item \textbf{Ferramentas:} Jest como \textit{test runner} e \textit{framework} de asserção; React Testing Library para testes de \textit{ViewModels} no frontend; Playwright para testes \textit{end-to-end}; \textit{Mocking} manual via \texttt{jest.mock} para isolar dependências externas (Supabase, PDFKit, \texttt{fetch}).
    \item \textbf{Execução em CI/CD:} A suíte completa é executada automaticamente pelo Render a cada \textit{push} na branch \texttt{main}. O pipeline valida que todos os testes passem antes da implantação.
    \item \textbf{Meta de cobertura:} O projeto define cobertura mínima de 75\% em \textit{Statements}, \textit{Functions} e \textit{Lines}, configurada no Jest e validada no pipeline de CI. As camadas de domínio e casos de uso no backend têm meta informal de 100\%.
    \item \textbf{Massa de dados:} Os testes de repositório utilizam o cliente Supabase \textit{mockado}, dispensando banco de dados real. Para testes das entidades de domínio, os dados são construídos em memória (\textit{in-memory}), garantindo isolamento total entre cenários.
\end{itemize}

\paragraph{Casos de Teste}

Cada caso de teste é identificado por um código único (CT\textit{nnn}) e descreve o cenário avaliado juntamente com o resultado esperado. O Quadro~\ref{qua:plano-testes} apresenta o plano completo, organizado por módulo e tipo de teste.

\renewcommand{\arraystretch}{0.85}
\small
\begingroup
\renewcommand{\tablename}{Quadro}
\renewcommand{\LTcaptype}{quadro}
\begin{longtable}{p{1.2cm} p{4.2cm} p{7.5cm} p{2.6cm}}
\caption{Plano de testes do sistema}
\label{qua:plano-testes}\\
\toprule
\textbf{Caso} & \textbf{Descrição} & \textbf{Resultados Esperados} & \textbf{Status} \\
\midrule
\endfirsthead
\toprule
\textbf{Caso} & \textbf{Descrição} & \textbf{Resultados Esperados} & \textbf{Status} \\
\midrule
\endhead
\bottomrule
\endfoot
\bottomrule
\endlastfoot

\multicolumn{4}{l}{\textbf{Módulo de Templates de Certificados (\ref{rf:upload-templates}--\ref{rf:duplicar-templates})}} \\
\midrule

CT001 & Criação de template & Template é criado com nome, descrição e \textit{layout} mesclado aos valores padrão. & Executado \\
CT002 & Listagem de templates & \textit{Endpoint} retorna lista de templates cadastrados. & Executado \\
CT003 & Obtenção de template por ID & \textit{Endpoint} retorna template específico; lança erro \texttt{TemplateNotFoundError} se não existir. & Executado \\
CT004 & Edição de template & Nome, descrição e configurações de \textit{layout} são alterados e persistem. & Executado \\
CT005 & Exclusão de template & Template é removido após confirmação do usuário; \textit{endpoint} \texttt{delete} é chamado. & Executado \\
CT006 & Personalização de \textit{layout} & Configurações parciais são mescladas com os valores padrão (\texttt{DEFAULT\_LAYOUT}). & Executado \\

\midrule
\multicolumn{4}{l}{\textbf{Módulo de Geração de Certificados (\ref{rf:gerar-pdf}--\ref{rf:download-lote})}} \\
\midrule

CT007 & Emissão manual de certificado & Selecionar template, preencher dados e confirmar gera certificado com código único. & Executado \\
CT008 & Geração de PDF & PDF gerado preserva layout, cores e fontes do template; resultado é um \textit{Buffer} binário. & Executado \\
CT009 & Geração de PDF com cores personalizadas & PDF é gerado com \texttt{backgroundColor}, \texttt{primaryColor} e \texttt{borderColor} personalizados. & Executado \\
CT010 & Geração de PDF sem borda & PDF é gerado corretamente com \texttt{showBorder = false}. & Executado \\
CT011 & Erro na geração de PDF & Falha no PDFKit rejeita a \textit{Promise} com mensagem de erro. & Executado \\
CT012 & Obtenção de certificado por ID & \textit{Endpoint} retorna certificado específico; lança erro \texttt{CertificateNotFoundError} se não existir. & Executado \\
CT013 & Listagem de certificados & \textit{Endpoint} retorna lista de certificados cadastrados. & Executado \\
CT014 & Atualização de certificado & Nome, CPF, carga horária, curso e \texttt{templateId} são alterados individual ou simultaneamente. & Executado \\
CT015 & Atualização de validade & Data de validade é atualizada; data passada é rejeitada. & Executado \\
CT016 & Validação de dados de entrada & Nome vazio, CPF inválido, carga horária negativa e \texttt{templateId} não-UUID são rejeitados. & Executado \\
CT017 & Validação de CPF & CPF com dígitos verificadores incorretos ou todos iguais é rejeitado. & Executado \\
CT018 & Formatação de CPF & CPF é automaticamente formatado (XXX.XXX.XXX-XX) durante a digitação. & Executado \\
CT019 & Emissão com template inexistente & \textit{Endpoint} lança \texttt{TemplateNotFoundError} se o \texttt{templateId} não existir. & Executado \\

\midrule
\multicolumn{4}{l}{\textbf{Módulo de Validação Pública e Antifraude (\ref{rf:portal-validacao}--\ref{rf:exibir-validacao})}} \\
\midrule

CT020 & Validação via código único & Código válido retorna certificado com status autêntico e válido. & Executado \\
CT021 & Certificado expirado & Certificado com validade vencida é retornado como autêntico, porém inválido. & Executado \\
CT022 & Código não encontrado & Código inexistente é rejeitado e a validação retorna certificado não localizado. & Executado \\
CT023 & \textit{Case insensitive} na consulta & Código digitado em minúsculo é convertido para maiúsculo antes da consulta no banco. & Executado \\
CT024 & Remoção de máscara do CPF & CPF com pontuação tem os caracteres não numéricos removidos antes da consulta. & Executado \\
CT025 & Tratamento de erro na validação & Erros retornados pela API são exibidos ao usuário e o estado de carregamento é gerenciado corretamente. & Executado \\

\midrule
\multicolumn{4}{l}{\textbf{Módulo de \textit{Dashboard} e Gestão (\ref{rf:lista-certificados}--\ref{rf:estatisticas-emissao})}} \\
\midrule

CT026 & Listagem de certificados & Tabela exibe certificados carregados da API com gerenciamento do estado de carregamento. & Executado \\
CT027 & Filtro por nome, CPF e código & Digitar texto no campo de busca filtra a lista em tempo real. & Executado \\
CT028 & Cópia do código de verificação & Botão copia o código para a área de transferência. & Executado \\
CT029 & Atualização de validade na interface & Usuário altera data de validade; a requisição é enviada ao servidor e a lista é recarregada. & Executado \\
CT030 & Estatísticas do \textit{dashboard} & Painel exibe total de certificados, ativos, templates e envios. & Executado \\
CT031 & Cálculo de certificados ativos & Certificados expirados são excluídos da contagem de ativos. & Executado \\

\midrule
\multicolumn{3}{l}{\textbf{Entidades de Domínio}} \\
\midrule

CT032 & Criação de certificado válido & Entidade \texttt{Certificate} é criada com todos os campos obrigatórios. & Executado \\
CT033 & Validações da entidade Certificate & Nome vazio, CPF inválido, curso vazio, carga horária zero/negativa e \texttt{templateId} vazio disparam erro de domínio. & Executado \\
CT034 & Data de validade futura & Certificado com validade futura retorna \texttt{isValid = true}. & Executado \\
CT035 & Data de validade passada & Certificado com validade passada retorna \texttt{isValid = false}. & Executado \\
CT036 & Criação de template válido & Entidade \texttt{Template} é criada com \textit{layout} mesclado aos valores \textit{default}. & Executado \\
CT037 & Validações da entidade Template & Nome vazio ou com apenas espaços dispara erro de domínio. & Executado \\
CT038 & Geração de código de verificação & Código gerado possui exatamente 8 caracteres do \textit{charset} definido. & Executado \\
CT039 & Comparação de códigos & Método \texttt{equals()} compara dois códigos corretamente; \textit{case insensitive}. & Executado \\
CT040 & Erros de domínio & Classes de erro (\texttt{CertificateNotFoundError}, \texttt{InvalidCPFError}, etc.) possuem mensagens e \textit{status} HTTP adequados. & Executado \\
CT041 & Restauração a partir de persistência & Dados vindos do banco são convertidos corretamente para instâncias das entidades. & Executado \\

\midrule
\multicolumn{3}{l}{\textbf{Camada de Infraestrutura}} \\
\midrule

CT042 & Repositório de certificados — consultas & Operações \texttt{findById}, \texttt{findByCPFAndCode} e \texttt{findAll} retornam certificados ou \texttt{null}. & Executado \\
CT043 & Repositório de certificados — escrita & Operações \texttt{create}, \texttt{update} e \texttt{delete} persistem no Supabase e tratam erros. & Executado \\
CT044 & Repositório de templates — consultas & Operações \texttt{findById} e \texttt{findAll} retornam templates ou \texttt{null} com \textit{fallback} de valores. & Executado \\
CT045 & Repositório de templates — escrita & Operações \texttt{create}, \texttt{update} e \texttt{delete} persistem no Supabase e tratam erros. & Executado \\
CT046 & Geração de PDF com PDFKit & PDF é gerado como \textit{Buffer}; cores e bordas são aplicadas; erro no PDFKit é tratado. & Executado \\
CT047 & Tratamento de erros da API & Erros de domínio (404, 400), erros Zod (400) e erros inesperados (500) são mapeados para \textit{status code} HTTP corretos. & Executado \\

\midrule
\multicolumn{3}{l}{\textbf{Utilitários (\textit{Frontend})}} \\
\midrule

CT048 & Chamadas HTTP da API & Métodos \textit{GET}, \textit{POST}, \textit{PUT}, \textit{DELETE} encapsulam \texttt{fetch} com tratamento de erro e \textit{parsing} JSON. & Executado \\
CT049 & Serviço de envios (\textit{localStorage}) & Operações de listagem e adição de registros de envio persistem no \textit{localStorage} com tratamento de erro. & Executado \\
CT050 & Detecção de \textit{mobile} & Hook \texttt{useIsMobile} retorna \texttt{true} para larguras menores que 768 px; \textit{listener} é limpo no desmonte. & Executado \\
CT051 & Função \texttt{isExpired} & Datas passadas retornam \texttt{true}; datas futuras e presentes retornam \texttt{false}. & Executado \\

\midrule
\multicolumn{3}{l}{\textbf{Testes Não-Funcionais}} \\
\midrule

CT052 & Carga — Geração simultânea de PDFs & Múltiplas requisições concorrentes de geração de PDF não causam falhas em cascata; o sistema permanece responsivo. & Parcial \\
CT053 & Segurança — Acesso não autenticado & Endpoints de criação, edição e exclusão de templates e certificados rejeitam requisições sem token de autenticação (status 401). & Executado \\
CT054 & Segurança — Força bruta na validação & Requisições repetidas ao endpoint de validação são limitadas (\textit{rate limiting}); o sistema bloqueia temporariamente após múltiplas tentativas inválidas. & Executado \\
CT055 & Acessibilidade — Navegação por teclado & Todos os campos do formulário de validação pública são acessíveis via tecla \texttt{Tab}; o botão de busca pode ser acionado via \texttt{Enter}. & Parcial \\

\midrule
\multicolumn{3}{l}{\textbf{Testes End-to-End (E2E)}} \\
\midrule

CT056 & Fluxo completo de emissão, validação e exclusão & Usuário cria template, emite certificado, valida o código no portal público e exclui o certificado com confirmação; todas as etapas são concluídas com sucesso. (Automatizado com Playwright.) & Executado \\
CT057 & Validação via \textit{QR Code} & \textit{QR Code} gerado no PDF, quando escaneado, redireciona para a URL de validação com o código correto como parâmetro. & Planejado \\

\midrule
\multicolumn{3}{l}{\textbf{Cenários de Resiliência}} \\
\midrule

CT058 & Indisponibilidade do Supabase & Cliente Supabase retorna erro de conexão; o sistema exibe mensagem amigável ao usuário sem interromper as demais funcionalidades. & Planejado \\
CT059 & Concorrência — Edição simultânea & Dois administradores editam o mesmo template ao mesmo tempo; o último salvamento prevalece e o conflito é registrado em log. & Executado \\
CT060 & \textit{Timeout} na geração de PDF & Geração de PDF excede o tempo limite; o sistema cancela a operação e notifica o usuário sobre a falha. & Planejado \\

\bottomrule
\end{longtable}
\par\vspace{4pt}
\endgroup
\noindent\small\textit{Fonte: Elaborado pelo Autor (2026).}
\normalsize
\renewcommand{\arraystretch}{1.0}

Ressalta-se que os casos não-funcionais CT052 a CT055 foram todos ao menos parcialmente verificados. O CT052 (carga) possui teste automatizado no backend com 10 chamadas concorrentes em menos de 5 segundos, porém sem variação de carga ou teste de estresse --- daí sua classificação como \textbf{Parcial}. O CT053 (autenticação) e o CT054 (\textit{rate limiting}) foram integralmente executados com testes automatizados no \textit{middleware}. O CT055 (acessibilidade) conta com teste automatizado de navegação Tab no frontend (\texttt{ValidityForm.test.tsx}), mas sem auditoria com ferramentas como \texttt{axe-core}, sendo classificado como \textbf{Parcial}.

Quanto ao CT057 (validação via \textit{QR Code}), sua marcação como \textbf{Planejado} deve-se ao fato de que, embora a geração do \textit{QR Code} no PDF já esteja implementada no \textit{backend} (pacote \texttt{qrcode} integrado ao \texttt{PDFKitService.ts}), não há um teste automatizado que valide que o código escaneado redireciona para a URL de verificação correta. A automação dessa verificação exigiria leitura óptica do \textit{QR Code} ou validação visual, o que está fora do escopo atual da suíte de testes E2E. A funcionalidade foi validada manualmente durante o desenvolvimento e permanece como melhoria para versões futuras.
